% ============================================================
\section{Introduction}
\label{sec:introduction}
% ============================================================

\subsection{Background and Motivation}
\label{subsec:background}

The dominant paradigm in digital advertising has, for two decades, relied on a
fundamental assumption: the user interacts with a visual interface---a web page,
a search results listing, a social media feed---where rectangular display regions
can be allocated to sponsored content. Banner advertisements, programmatic
display auctions, and search-engine sponsored links all exploit this property.
The advertiser purchases screen real estate; the publisher renders a creative
unit; the user either attends to it or does not. This model, refined through
successive generations of behavioural targeting and real-time bidding, generates
hundreds of billions of dollars in annual revenue worldwide and underpins the
economic viability of most free-to-use internet services.

Conversational AI interfaces disrupt this assumption at its foundation. Systems
such as ChatGPT, Claude, and Gemini present information as a single stream of
natural-language text. There are no sidebars, no interstitial slots, and no
visual regions in which a banner can be placed without breaking the
conversational contract between the system and its user. At the same time, these
interfaces are absorbing an increasing share of the queries that users
previously directed to search engines---product research, comparison shopping,
recommendation requests---precisely the query types that carry the highest
commercial intent and, consequently, the highest advertising value. The shift is
not merely a change of medium; it is a structural transformation in how users
express intent, receive information, and make purchasing decisions.

The advertising industry therefore faces a concrete technical problem: how to
monetise conversational AI interactions without degrading the user experience
that makes them valuable in the first place. Appending a block of sponsored
links to the end of a chat response replicates the search-ad format but
sacrifices the natural feel of the conversation. Inserting interstitial
advertisements between turns interrupts the dialogue flow and erodes user trust.
Neither approach exploits the distinctive capability of a large language model
(LLM), namely its ability to generate contextually appropriate, fluent
natural-language text that integrates information from multiple sources---including
sponsored product information---into a single coherent response.

This observation motivates the present work. If the LLM can understand the
user's intent, match it against an inventory of relevant products, and weave a
sponsored recommendation naturally into an otherwise helpful response, then the
advertisement ceases to be an interruption and becomes, instead, part of the
answer. The result is an advertising format that is native to the conversational
medium: contextually relevant, linguistically integrated, and measurable not
merely through impressions and clicks but through the user's subsequent
conversational engagement with the recommended product.

\subsection{Project Objectives}
\label{subsec:objectives}

This project designs, implements, and evaluates an end-to-end LLM-powered
conversational advertising platform---referred to throughout this report as
\textit{Ad LLM as a Service}---that delivers inline sponsored product
recommendations within AI chat responses. The platform satisfies the following
objectives:

\begin{enumerate}[label=(\roman*), leftmargin=2.5em]
  \item \textbf{Intent-aware context extraction.} A lightweight LLM call
    extracts structured context from each user message and the preceding
    conversation history, producing a classification of intent, topical
    relevance, sentiment, purchase signal strength, and---critically---a
    sensitivity-aware ad receptivity score that gates whether any advertisement
    should be shown at all.

  \item \textbf{Semantic ad matching.} The extracted context is embedded into a
    vector representation and matched against a pre-embedded ad inventory via
    cosine similarity search over \texttt{pgvector}, followed by a business-rule
    re-ranking stage that incorporates bid price, budget constraints, and
    per-session frequency caps.

  \item \textbf{Natural ad integration with mandatory disclosure.} The main
    response-generation LLM receives the matched advertisement as part of its
    prompt context and is instructed to weave the product mention into a helpful
    answer if and only if the product is genuinely relevant. Every sponsored
    mention carries a \texttt{[Sponsored]} disclosure tag, ensuring compliance
    with Federal Trade Commission (FTC) endorsement guidelines.

  \item \textbf{Conversational engagement measurement.} Beyond conventional
    impression and click-through tracking, the system detects whether a user's
    follow-up message engages with the advertised product---a signal that is
    unique to conversational interfaces and substantially stronger than a
    passive impression.

  \item \textbf{Ethical safeguards.} Conversations classified as sensitive---those
    involving health crises, emotional distress, financial hardship, or
    grief---are automatically excluded from ad injection, regardless of
    commercial relevance.

  \item \textbf{Functional MVP.} The deliverable is a working prototype that
    demonstrates the complete pipeline from user message to ad-augmented
    response, accompanied by an analytics dashboard reporting impressions,
    click-through rates, engagement events, and estimated revenue.
\end{enumerate}

\subsection{Scope and Constraints}
\label{subsec:scope}

The platform is scoped as a minimum viable product (MVP) and is subject to
several deliberate constraints that bound the design space.

The architecture is a single-service deployment: one FastAPI orchestrator
coordinates all internal modules, a single Redis instance manages session state,
and a single PostgreSQL database (extended with \texttt{pgvector}) stores both
the ad inventory and the event log. This monolithic design is appropriate for the
current scale---tens of thousands of ads, moderate request concurrency---and
avoids the operational complexity of a distributed microservice architecture.

The ad inventory is pre-loaded via a bulk ingestion script; no advertiser-facing
self-service portal is included in the MVP. Campaign management, budget pacing,
and creative upload are deferred to subsequent iterations.

Two external dependencies impose hard constraints on latency and cost. Every
request requires at minimum two LLM API calls---one lightweight call for context
extraction and one full-model call for response synthesis---plus one embedding
API call for the query vector. The total end-to-end latency budget is three
seconds, inclusive of all network round-trips and internal processing. The
per-request cost is dominated by these API calls, and the re-ranking stage uses
heuristic scoring weights (a weighted combination of semantic similarity, normalised
bid price, and purchase signal) rather than a learned model, since no
click-through data is available at launch.

The engagement classifier, which determines whether a user's follow-up message
references the advertised product, is initialised as a keyword-and-entity
matching module. A lightweight LLM-based classifier is planned as an upgrade
once sufficient labelled engagement data has been collected.

\subsection{Report Organisation}
\label{subsec:report-org}

The remainder of this report is organised as follows.
Section~\ref{sec:literature-review} surveys the relevant literature across four
areas: the evolution of digital advertising formats, the capabilities of large
language models for controlled generation, retrieval-augmented generation as a
paradigm for context-aware information retrieval, and the emerging body of work
on advertising within conversational interfaces. Section~\ref{sec:problem-statement}
formulates the problem precisely, identifying the technical challenges of
real-time intent extraction, low-latency semantic ad matching, natural ad
integration, and engagement measurement, together with the ethical and regulatory
constraints that any solution must satisfy.
Section~\ref{sec:solution-design} presents the system architecture and the
design of each component: the FastAPI orchestrator, the context extraction
module, the Redis-backed session manager, the \texttt{pgvector}-based ad
matching pipeline, the prompt-engineered response synthesis module, and the
analytics and click-tracking subsystem. Section~\ref{sec:implementation}
describes the implementation in detail, including workstream responsibilities,
the database schema, API contracts, and the quality assurance framework.
Section~\ref{sec:illustrative-examples} walks through five end-to-end examples
that exercise representative paths through the pipeline, including successful ad
injection, sensitive-conversation suppression, irrelevant-ad rejection,
engagement follow-up detection, and multi-turn context accumulation.
Section~\ref{sec:conclusions} summarises the contributions, discusses
limitations, and outlines directions for future work.